\hspace{3mm}

\centerline{\bf \Huge Первая тренировочная Олимпиада Эйлера}

\hspace{3mm}

 \begin{flushright}
     \scriptsize{Любимое дело придает нашей жизни смысл.\\Кенни Аккерман}
 \end{flushright}
 
\problem{1} Машина Герман бежал по трассе $P216$ из Элисты в Астрахань со скоростью $70$ км/ч и попал под песчанную бурю, которая двигалась вдоль трассы в ту же сторону, что и Машина Герман (но медленнее его). Встречная Машина Эрдем двигалась из Астрахани в Элисту по той же трассе и с той же скоростью, что и Машина Герман, но пробыл под бурей в три раза меньше, чем Герман. С какой сторостью двигалась буря?

\textit{Примечание:} песчанная буря это некоторый объект с определенной длинной, который двигается вдоль трассы.

\hspace{3mm}

\problem{2}За круглым столом в очередной раз собрались физматы и гумы, причем каждых было по 100 человек. Для каждого физмата оказалось, что среди $12$ следующих за ним по часовой стрелке людей физматов столько же, сколько и среди $12$ предыдущих людей. Могло ли так оказаться, что для гума количество физматов среди следующих $12$ и предыдущих $12$ отличаются ровно на $1$?

\hspace{3mm}

\problem{3}Дано $6$ простых чисел $p_1, p_2, \dotsc, p_6$ таких, что $p_{i+1} = 2p_i + 1$. Может ли сумма попарных произведений этих чисел быть степенью простого числа?

\hspace{3mm}

\problem{4} Точка $M -$ середина стороны $AC$ треугольника $ABC.$ На сторонах $AB$ и $BC$ выбраны точки $X$ и $Y$ соответственно так, что $\angle XMY = 90^{\circ}$. Докажите, что \\

\centerline{$AM -MX > CY - YX$}

\hspace{3mm}

\problem{5}В художественном клубе ЛмШ осталось $k$ красок. Расул Владимирович заметил, что при любом разделении участников клуба на две группы (непустые) существует такой набор красок, что каждый член одной группы использовал чётное количество красок из этого набора, чтобы написать картину, а каждый член другой группы нечётное количество красок. Какое наибольшее количество участников в клубе ЛмШ?